\section{Gauge Group Derivation: Standard Model Closure from MAR plus Explicit Gauge Premises}\label{gauge-group-derivation-standard-model-closure-under-extended-mar}

\begin{quote}
This part starts from the five OPH axioms, the canonical technical premises, the fixed-cutoff realized presentation, and the compact-gauge reconstruction theorem, and then applies the admissibility analysis that closes the realized Standard Model branch.

\textbf{Overview.} Once Tannaka/DR reconstruction~\cite{tannaka38,krein49,dr89,dr90} is applied to the refinement-stable edge-sector colimit, the gauge-selection hypotheses pick out the Standard Model connected Lie gauge-sector factors once the admissible class includes one connected abelian charge factor. On the realized admissible branch, the CP/UV window and MAR fix the generation count, Witten's \(\mathrm{SU}(2)\) anomaly parity~\cite{witten82} then fixes the color count, and the realized hypercharge lattice fixes the global gauge quotient:

\[G_{\mathrm{phys}} = \frac{SU(3) \times SU(2) \times U(1)}{\mathbb{Z}_6}, \qquad N_g = 3, \qquad N_c = 3.\]

Within the positive-dimensional connected Lie admissible class, no other gauge-sector package, no other color multiplicity, and no other generation count satisfies all admissibility conditions while minimizing the complexity vector \(C(\mathfrak S)\).
\end{quote}

\begin{center}\rule{0.5\linewidth}{0.5pt}\end{center}

\subsection{Gauge-Selection Setup}\label{1-premise-summary}

The gauge-reconstruction theorem, the central transport-obstruction clause, and the fixed-cutoff realized presentation are assumed throughout. The analysis below identifies which compact reconstructed group survives the admissibility filters and Axiom 5 (MAR).

\begin{center}\rule{0.5\linewidth}{0.5pt}\end{center}

\subsection{Admissibility Conditions}\label{2-admissibility-conditions}

An \textbf{OPH-realizable gauge sector package} \(\mathfrak S = (G, \text{matter content})\) consists of a compact gauge group \(G\) reconstructed via Tannaka-Krein~\cite{tannaka38,krein49,dr89,dr90} from the refinement-stable edge-sector category (Theorem 6.1) together with its associated chiral matter spectrum. A sector package \(\mathfrak S\) is \textbf{admissible} if it satisfies all of the following conditions:

\subsubsection{(i) Loop-coherent / transportable}\label{i-loop-coherent--transportable}

The central obstruction class vanishes: \([z] = 0\). Equivalently, edge charges are DHR-transportable, i.e., they can be moved between patches without affecting fusion rules. This is the content of the explicit premise \([z] = 0\) (Proposition 6.1a).

\emph{Rationale:} Without transportability, the reconstructed gauge symmetry is only a local labeling, not a global symmetry of the theory.

\subsubsection{(ii) Anomaly-free}\label{ii-anomaly-free}

The gauge sector is free of all perturbative anomalies (ABJ anomaly cancellation) and all global anomalies (Witten SU(2) anomaly~\cite{witten82}, Dai-Freed anomalies). Concretely:

\begin{itemize}
\tightlist
\item
  \(\mathrm{tr}[T_a \{T_b, T_c\}] = 0\) for all gauge generators (perturbative),
\item
  \(\pi_4(G)\) anomaly vanishes (Witten\textquotesingle s global SU(2) anomaly),
\item
  Mixed gravitational-gauge anomalies cancel.
\end{itemize}

\emph{Rationale:} Anomalous gauge theories are inconsistent at the quantum level.

\subsubsection{(iii) Refinement-stable with light chiral matter}\label{iii-refinement-stable-with-light-chiral-matter}

The MaxEnt/refinement-stable state supports light charged fermions without fine-tuning. This requires the matter content to be \emph{chiral} , i.e., left-handed and right-handed fermions carry different gauge representations , so that vector-like mass terms are forbidden by gauge symmetry.

\emph{Rationale:} The refinement-stability lemma in the main OPH trunk shows that refinement stability forbids unprotected relevant operators. A gauge-invariant Dirac mass term is relevant; if both chiralities are in conjugate representations, the mass term is allowed and drives the fermions to the cutoff. The corresponding chirality corollary then selects chiral content as the natural refinement-stable option.

\subsubsection{(iv) Single-Higgs Yukawa-completable}\label{iv-single-higgs-yukawa-completable}

The chiral matter content can acquire mass through Yukawa couplings to a single scalar doublet \(H\) carrying some charge under one connected abelian factor, after electroweak symmetry breaking, without requiring additional scalar multiplets.

\emph{Rationale:} This is the minimal mechanism for generating fermion masses consistent with electroweak symmetry breaking. The admissibility premise only requires one connected abelian charge direction and a single Higgs doublet; the numerical hypercharge assignments are derived later from anomaly cancellation and Yukawa invariance. Additional Higgs multiplets would increase the scalar capacity without physical necessity.

\subsubsection{(v) Intrinsically CP-capable}\label{v-intrinsically-cp-capable}

The sector supports intrinsic CP violation , i.e., CP-violating phases that cannot be removed by field redefinitions. For a CKM-type mixing matrix with \(N_g\) generations, the number of physical CP phases is \((N_g - 1)(N_g - 2)/2\).

\emph{Rationale:} Observed CP violation in the kaon and B-meson systems requires intrinsic CP-violating phases. A sector that cannot accommodate them is empirically excluded.

\subsubsection{(vi) Weak-sector UV-completable}\label{vi-weak-sector-uv-completable}

The weak gauge sector (the factor carrying the pseudoreal doublet) is asymptotically free at one loop, ensuring that it is UV-completable within the OPH framework rather than requiring a Landau pole cutoff below the screen scale.

\emph{Rationale:} The OPH screen operates at the Planck scale. A gauge factor with a sub-Planckian Landau pole would be inconsistent with the screen construction.

\begin{center}\rule{0.5\linewidth}{0.5pt}\end{center}

\subsection{Axiom 5 (MAR): Formal Statement Used Here}\label{3-selection-axiom-mar}

\textbf{Axiom 5 (MAR): Minimal Admissible Realization.} Among all OPH-realizable sector packages \(\mathfrak S\) that satisfy admissibility conditions (i)--(vi), the realized low-energy package is the \textbf{lexicographically minimal} one under the complexity vector

\[C(\mathfrak S) = \bigl(\chi_{\mathrm{cpl}},\; N_{\mathrm{nonab}},\; N_c,\; N_g\bigr),\]

where:

\begin{itemize}
\tightlist
\item
  \(\chi_{\mathrm{cpl}}\) is the \textbf{coupled edge capacity}: the dimension of the minimal unitary carrier containing a common irreducible block on which the admissible pseudoreal and complex nonabelian charge types both act nontrivially,
\item
  \(N_{\mathrm{nonab}}\) is the \textbf{number of nonabelian simple factors} in \(G\),
\item
  \(N_c\) is the \textbf{rank of the color factor} (the dimension of the fundamental representation of the complex nonabelian factor),
\item
  \(N_g\) is the \textbf{number of chiral generations}.
\end{itemize}

This subsection fixes the exact formal statement of Axiom 5 used by the gauge proof.

Lexicographic minimality means: first minimize \(\chi_{\mathrm{cpl}}\); among ties, minimize \(N_{\mathrm{nonab}}\); among further ties, minimize \(N_c\); among final ties, minimize \(N_g\). The object minimized here is the full sector package \(\mathfrak S\), not the bare tensor category alone. MAR is therefore an explicit selector on the admissible class, not a theorem derived from the earlier axioms.

\textbf{Key properties of MAR:}

\begin{enumerate}
\def\labelenumi{\arabic{enumi}.}
\item
  \emph{It is a selection rule, not a dynamics.} MAR acts on an independently defined admissible class; it does not modify the OPH equations of motion.
\item
  \emph{It is not "minimalism in general."} MAR only selects among sectors that pass all six admissibility filters. Plain minimality (smallest \(G\)) would yield \(U(1)\) or \(SU(2)\), which fail conditions (iii)--(v). The admissibility conditions do the heavy lifting; MAR breaks the remaining degeneracy.
\item
  \emph{It is powerful.} Product gauge structure and the specific factors \(SU(3) \times SU(2) \times U(1)\) follow from MAR applied to the admissible class; on the realized branch the CP/UV window and MAR then give \(N_g = 3\), and Witten parity with that realized generation count gives \(N_c = 3\).
\end{enumerate}

\textbf{Remark (coupled vs. faithful).} The first MAR coordinate is the coupled carrier dimension \(\chi_{\mathrm{cpl}}\), not the abstract minimal faithful representation dimension. For \(G_{\mathrm{phys}} \cong S(U(3) \times U(2))\), the block-diagonal action on \(\mathbb{C}^3 \oplus \mathbb{C}^2\) is faithful of dimension \(5\), but it places the complex and pseudoreal charge types on disconnected invariant subspaces. MAR minimizes the coupled carrier dimension \(\chi_{\mathrm{cpl}}\), for which the minimal SM carrier is \(\mathbb{C}^3 \otimes \mathbb{C}^2\) of dimension \(6\).

\begin{center}\rule{0.5\linewidth}{0.5pt}\end{center}

\subsection{Main Theorem}\label{4-main-theorem}

\textbf{Theorem (Standard Model closure from the gauge-selection hypotheses).} Under the gauge-reconstruction hypotheses of Theorem 6.1, the central transport condition \([z]=0\), admissibility conditions (i)--(vi), and Axiom 5 (MAR):

\[G_{\mathrm{phys}} = \frac{SU(3) \times SU(2) \times U(1)}{\mathbb{Z}_6}, \qquad N_g = 3, \qquad N_c = 3.\]

The realized counting chain is sequential rather than joint: the CP/UV admissibility window and MAR first select \(N_g = 3\), and Witten's \(\mathrm{SU}(2)\) anomaly parity then fixes \(N_c = 3\).

\emph{The proof is given in \S{}5--\S{}7 below.}

\begin{center}\rule{0.5\linewidth}{0.5pt}\end{center}

\subsection{Proof: Gauge Group}\label{5-proof-gauge-group}

The proof proceeds in five steps. Steps 1--4 determine the connected gauge group. Step 5 fixes the global structure.

\subsubsection{Step 1: Some compact group exists}\label{step-1-some-compact-group-exists}

\textbf{From:} the canonical OPH package together with the fixed-cutoff realized presentation.

At each cutoff one has a finite tensor category of transportable edge charges. The local MaxEnt/refinement package controls the realized state branch across cutoffs, but it does not by itself prove that transportable sectors persist nontrivially in the refinement limit. The reconstruction therefore uses the transportable directed-colimit premise \(T6\): assume a directed system of transportable edge-sector categories with objectwise finite-dimensional fibers, and write its colimit as \(\mathsf{Sect}_\infty\). By Theorem 6.1 (Tannaka/DR reconstruction~\cite{tannaka38,krein49,dr89,dr90}), there exists a compact group \(G\), unique up to isomorphism, with \(\mathsf{Sect}_\infty \simeq \mathrm{Rep}(G)\).

The premise \([z] = 0\) ensures that this reconstruction is global: charges are DHR-transportable, so the compact group acts as a genuine gauge symmetry, not merely a local labeling (Proposition 6.1a). \(\square\)

\subsubsection{Step 2: The gauge group must be a product}\label{step-2-the-gauge-group-must-be-a-product}

\textbf{From:} MAR (minimizing \(\chi_{\mathrm{cpl}}\)) + admissibility (iii).

\textbf{Lemma 5.1.} Any admissible sector must contain at least two genuinely different nonabelian charge types: one pseudoreal and one complex.

\emph{Proof.} Admissibility condition (iii) requires light chiral matter. By the corresponding chirality corollary in the main OPH trunk, refinement stability forces the matter content to be chiral. To support chiral fermions that can form gauge-invariant Yukawa couplings with a single Higgs doublet (condition iv), the gauge group must admit both:

\begin{itemize}
\tightlist
\item
  A \textbf{pseudoreal} fundamental representation (for the weak doublet structure , this allows left-handed doublets and right-handed singlets to have different gauge quantum numbers), and
\item
  A \textbf{complex} fundamental representation (for the color structure , this allows quarks to carry a charge that distinguishes particle from antiparticle).
\end{itemize}

A single nonabelian factor cannot simultaneously provide both types: a group is either pseudoreal at a given dimension or complex, not both. \(\square\)

\textbf{Lemma 5.2.} The minimal coupled carrier is \(V = \mathbb{C}^3 \otimes \mathbb{C}^2\), with \(\chi_{\mathrm{cpl}} = 6\).

\emph{Proof.}

\begin{itemize}
\tightlist
\item
  The smallest faithful pseudoreal irreducible representation of any positive-dimensional compact connected Lie group is the fundamental \textbf{2} of \(SU(2)\) (the only 2D pseudoreal irrep in the connected Lie setting used here).
\item
  No irreducible 2D unitary representation of a positive-dimensional compact connected Lie group can realize the required intrinsically complex nonabelian role: the connected derived subgroup of its image is a compact connected subgroup of \(SU(2)\), hence either trivial, a torus, or all of \(SU(2)\), so the nonabelian action is pseudoreal up to abelian twist. The smallest faithful irreducible intrinsically complex nonabelian representation is therefore the fundamental \textbf{3} of \(SU(3)\).
\item
  By definition of \(\chi_{\mathrm{cpl}}\), the admissible pseudoreal and complex charge types must occur on a common nonabelian block rather than on disconnected direct-sum summands. The minimal such block is the tensor product \(V = \mathbb{C}^3 \otimes \mathbb{C}^2\), giving \(\chi_{\mathrm{cpl}} = 6\).
\item
  The block-diagonal representation \(\mathbb{C}^3 \oplus \mathbb{C}^2\) of \(S(U(3) \times U(2))\) is faithful in the abstract group-theoretic sense, but it is not coupled: each nonabelian factor is trivial on one summand. It therefore does not enter the MAR selector.
\item
  Any other combination (e.g., using \(SU(4)\) instead of \(SU(3)\), or \(Sp(4)\) instead of \(SU(2)\)) gives \(\chi_{\mathrm{cpl}} > 6\), and is therefore disfavored by MAR. \(\square\)
\end{itemize}

\textbf{Proposition 5.3 (Product structure from minimal coupled carrier).} The minimal coupled carrier \(\mathbb{C}^3 \otimes \mathbb{C}^2\) forces a product gauge structure.

\emph{Proof.} On \(V = \mathbb{C}^3 \otimes \mathbb{C}^2\), the color factor acts on the first tensor factor and the weak factor acts on the second. These actions commute by the tensor product structure. Therefore the gauge group decomposes as a product \(G_{\mathrm{color}} \times G_{\mathrm{weak}} \times G_{\mathrm{abelian}}\) (up to finite quotient). A simple group acting irreducibly on a low-dimensional carrier would not provide the independent pseudoreal + complex factor structure required by Lemma 5.1; when those roles are realized on a common block, the minimal coupled carrier is the tensor product above.

Product structure is derived from the minimal coupled-carrier argument. \(\square\)

\textbf{Remark.} The representation-theoretic dimension comparisons in Steps 2--4 apply only to the positive-dimensional connected Lie image carrying the admissible nonabelian charges, equivalently to the identity component of the reconstructed group on that image. Finite or disconnected counterexamples are outside the theorem package used here.

\subsubsection{Step 3: The factors are SU(3), SU(2), and U(1)}\label{step-3-the-factors-are-su3-su2-and-u1}

\textbf{From:} The core gauge lemmas of the main OPH trunk, applied to the minimal coupled carrier together with the connected abelian admissibility premise.

\textbf{Lemma 5.4 (SU(2) from pseudoreal doublet).} The factor acting on \(\mathbb{C}^2\) as a faithful 2D pseudoreal representation is \(SU(2)\).

\emph{Proof.} Among compact groups with a faithful 2D pseudoreal unitary representation, \(SU(2)\) is the unique option. (The only compact Lie groups with a 2D faithful representation are subgroups of \(U(2)\). Among these, the pseudoreal condition \(V \cong \bar{V}\) via an antisymmetric bilinear form selects exactly \(SU(2)\).) This is the SU(2)-selection lemma of the main OPH trunk. \(\square\)

\textbf{Lemma 5.5 (SU(3) from complex triplet).} The factor acting on \(\mathbb{C}^3\) as a faithful irreducible complex representation is \(SU(3)\).

\emph{Proof.} Among compact groups with a faithful irreducible complex 3D representation, the connected component is \(SU(3)\). (The only compact simple Lie groups with a 3D fundamental representation are \(SU(3)\) and \(SO(3)\). But the fundamental of \(SO(3)\) is real, not complex. Therefore \(SU(3)\) is uniquely selected.) This is the SU(3)-selection lemma of the main OPH trunk. \(\square\)

\textbf{Lemma 5.6 (Single connected abelian factor).} Admissibility condition (iv) supplies one connected abelian charge factor acting nontrivially on the coupled carrier. Compact connected abelian Lie groups are tori, and the commutant calculation below shows that only a one-torus can act nontrivially on the minimal coupled carrier. Therefore the admissible connected abelian factor is \(U(1)\).

\emph{Proof.} This is the connected-abelian-factor lemma of the main OPH trunk together with the compact-Lie classification of connected abelian groups as tori. The one-dimensional characters form a discrete lattice \(\mathrm{Hom}(U(1)^n,U(1)) \cong \mathbb{Z}^n\); what matters here is the connected torus factor, not a continuous family of inequivalent 1D irreps. \(\square\)

\subsubsection{Step 4: Nothing else , the commutant argument}\label{step-4-nothing-else--the-commutant-argument}

\textbf{Proposition 5.7 (No additional factors).} No extra gauge factor , neither nonabelian nor connected abelian , can appear without increasing \(\chi_{\mathrm{cpl}}\) beyond 6.

\emph{Proof.} Consider the maximal compact subgroup of \(U(6)\) acting on \(V = \mathbb{C}^3 \otimes \mathbb{C}^2\) with commuting actions on each tensor factor:

\[\{g \in U(6) : g = g_3 \otimes g_2,\; g_3 \in U(3),\; g_2 \in U(2)\} \cong U(3) \times U(2) / U(1)_{\mathrm{diag}}.\]

The continuous symmetry group with commuting color and weak actions is therefore

\[S(U(3) \times U(2)) \cong \frac{SU(3) \times SU(2) \times U(1)}{\text{finite center}}.\]

Compute the commutant:

\begin{itemize}
\tightlist
\item
  The commutant of \(SU(3) \times SU(2)\) inside \(U(6)\) is exactly \(U(1)\).
\item
  Therefore any connected abelian factor acting nontrivially on the coupled carrier is necessarily this single \(U(1)\), and no extra continuous factor (neither \(U(1)'\) nor any nonabelian group) can appear without enlarging the carrier beyond \(\chi_{\mathrm{cpl}} = 6\).
\item
  Any additional nonabelian factor would have to act nontrivially on the same coupled block, increasing \(\chi_{\mathrm{cpl}}\).
\end{itemize}

MAR selects the minimal \(\chi_{\mathrm{cpl}}\), so extra factors are excluded. \(\square\)

\subsubsection{\texorpdfstring{Step 5: The global quotient is \(\mathbb{Z}_6\)}{Step 5: The global quotient is \textbackslash mathbb\{Z\}\_6}}\label{step-5-the-global-quotient-is-mathbbz_6}

\textbf{Proposition 5.8 (Z\textsubscript{6} quotient from hypercharge quantization).} If the realized matter spectrum has hypercharges quantized in sixths, then the subgroup of \(SU(3) \times SU(2) \times U(1)\) acting trivially on all physical states is exactly \(\mathbb{Z}_6\).

\emph{Proof.} The SM matter content with hypercharge assignments \((Q_L, u_R, d_R, L_L, e_R, H)\) has hypercharges \(Y \in \{1/6, 2/3, -1/3, -1/2, -1, 1/2\}\), all in \(\frac{1}{6}\mathbb{Z}\). The center of \(SU(3) \times SU(2) \times U(1)\) is \(\mathbb{Z}_3 \times \mathbb{Z}_2 \times U(1)\). The subgroup acting trivially on all realized representations is generated by

\[(\omega_3, -1, e^{i\pi/3}) \in SU(3) \times SU(2) \times U(1),\]

where \(\omega_3 = e^{2\pi i/3}\) is the \(\mathbb{Z}_3\) center of \(SU(3)\). This generates a \(\mathbb{Z}_6\) subgroup. Therefore

\[G_{\mathrm{phys}} = \frac{SU(3) \times SU(2) \times U(1)}{\mathbb{Z}_6}.\]

Equivalently, \(G_{\mathrm{phys}} \cong S(U(3) \times U(2))\).

The hypercharge quantization itself follows from anomaly cancellation (condition ii) with the minimal Yukawa-complete spectrum (condition iv). The hypercharge theorem in the main OPH trunk derives the assignments from these conditions for \(N_c = 3\), yielding precisely the sixth-integer lattice. This is the same proposition used there. \(\square\)

\subsubsection{Summary of Step 1--5}\label{summary-of-step-15}

\[\boxed{\text{gauge-selection hypotheses} \;\Longrightarrow\; G_{\mathrm{phys}} = \frac{SU(3) \times SU(2) \times U(1)}{\mathbb{Z}_6}}\]

\begin{center}\rule{0.5\linewidth}{0.5pt}\end{center}

\subsection{Proof: Number of Colors}\label{6-proof-number-of-colors}

\textbf{Theorem (N\_c = 3).} Under the gauge-selection hypotheses above, the rank of the color factor is \(N_c = 3\).

\textbf{Proof.} The gauge group derivation (\S{}5) fixes the color factor as \(SU(3)\), i.e. \(N_c = 3\), through the minimal coupled-carrier argument (Lemma 5.2). We give here the independent confirmation from anomaly constraints plus MAR.

\textbf{Step 1: Witten anomaly constrains N\_c once N\_g is fixed.}

With gauge structure \(SU(N_c) \times SU(2)_L \times U(1)_Y\), one left-handed quark doublet \(Q\) per color, and one left-handed lepton doublet \(L\) per generation, the total number of \(SU(2)\) doublets is

\[N_{\mathrm{doublets}} = N_g(N_c + 1).\]

Using \(N_g = 3\) from Section 7, Witten\textquotesingle s global \(SU(2)\) anomaly~\cite{witten82} requires the total number of \(SU(2)\) doublets to be even:

\[3(N_c + 1) \equiv 0 \pmod{2} \quad \Longrightarrow \quad N_c \text{ is odd}.\]

This alone allows \(N_c \in \{1, 3, 5, 7, \ldots\}\).

\textbf{Step 2: N\_c = 1 fails admissibility.}

For \(N_c = 1\): \(SU(1)\) is trivial (no color dynamics). The fundamental representation is 1-dimensional and real, not complex. This violates Lemma 5.1 (no genuinely complex nonabelian charge type) and condition (iii) (cannot support chiral quarks).

\textbf{Step 3: MAR selects N\_c = 3.}

Among the remaining candidates \(\{3, 5, 7, \ldots\}\), the SM-like quotient family has \(\chi_{\mathrm{cpl}} = 2N_c\) because the minimal coupled carrier is \(\mathbb{C}^{N_c} \otimes \mathbb{C}^2\). (The block-diagonal faithful representation \(\mathbb{C}^{N_c} \oplus \mathbb{C}^2\) has dimension \(N_c + 2\), but it is not coupled and is therefore irrelevant to MAR.) MAR\textquotesingle s complexity vector has \(N_c\) in the third slot, so lexicographic minimization selects \(N_c = 3\). \(\square\)

\begin{center}\rule{0.5\linewidth}{0.5pt}\end{center}

\subsection{Proof: Number of Generations}\label{7-proof-number-of-generations}

\textbf{Theorem (N\_g = 3).} Under the gauge-selection hypotheses above, the number of chiral generations is \(N_g = 3\).

\textbf{Proof.}

\textbf{Step 1: CP violation gives N\_g \ensuremath{\geq} 3.}

Admissibility condition (v) requires intrinsic CP violation. The number of physical CP-violating phases in an \(N_g \times N_g\) CKM matrix is

\[n_{\mathrm{CP}} = \frac{(N_g - 1)(N_g - 2)}{2}.\]

\begin{itemize}
\tightlist
\item
  For \(N_g = 1\): \(n_{\mathrm{CP}} = 0\). No CP violation possible. Excluded.
\item
  For \(N_g = 2\): \(n_{\mathrm{CP}} = 0\). No CP violation possible. Excluded.
\item
  For \(N_g = 3\): \(n_{\mathrm{CP}} = 1\). CP violation possible.
\end{itemize}

Therefore \(N_g \geq 3\).

\textbf{Step 2: Asymptotic freedom gives N\_g \ensuremath{\leq} 5.}

Admissibility condition (vi) requires \(SU(2)_L\) to be asymptotically free. With one complex Higgs doublet, the one-loop beta function coefficient is

\[b_{SU(2)} = \frac{22}{3} - \frac{1}{3}N_g(N_c + 1) - \frac{1}{6}.\]

Asymptotic freedom requires \(b_{SU(2)} > 0\):

\[N_g(N_c + 1) < \frac{43}{2}.\]

Since the color-type role forces \(N_c \ge 3\), one has \(N_c + 1 \ge 4\). Therefore \(4N_g < \frac{43}{2}\), so \(N_g \leq 5\).

Combining: \(3 \leq N_g \leq 5\).

\textbf{Step 3: MAR selects N\_g = 3.}

Among the admissible candidates \(\{3, 4, 5\}\), MAR\textquotesingle s complexity vector has \(N_g\) in the fourth (last) slot. Lexicographic minimization selects \(N_g = 3\). Feeding that value back into Section 6 completes the Witten-parity closure \(N_c = 3\). \(\square\)

\begin{center}\rule{0.5\linewidth}{0.5pt}\end{center}

\subsection{Corollaries}\label{8-corollaries}

\subsubsection{Corollary 1: No gauge-mediated proton decay}\label{corollary-1-no-gauge-mediated-proton-decay}

\textbf{Corollary.} The gauge group \(G_{\mathrm{phys}} = SU(3) \times SU(2) \times U(1)/\mathbb{Z}_6\) does not contain \(X\) or \(Y\) gauge bosons that mediate proton decay.

\emph{Proof.} The adjoint representation of the full connected gauge group contains only \((8,1,0)\oplus(1,3,0)\oplus(1,1,0)\); equivalently, the adjoint of the derived gauge group contains only \((8,1,0)\oplus(1,3,0)\). There are therefore no mixed \((3,2,\pm 5/6)\) generators. Proton decay via gauge bosons requires a simple unification group (e.g., \(SU(5)\), \(SO(10)\)) whose additional generators connect quarks to leptons. The MAR derivation forces the product-group structure above, so no such gauge bosons exist. \(\square\)

\subsubsection{Corollary 2: Uniqueness}\label{corollary-2-uniqueness}

\textbf{Corollary.} Within the positive-dimensional connected Lie admissible class with one connected abelian factor, the SM gauge group is the unique realization acting on the minimal coupled carrier with \(\chi_{\mathrm{cpl}} = 6\).

\emph{Proof.} By the commutant argument (Proposition 5.7), the group \(SU(3) \times SU(2) \times U(1)/\mathbb{Z}_6\) exhausts the continuous symmetry of the minimal coupled carrier \(\mathbb{C}^3 \otimes \mathbb{C}^2\). Any alternative group either:

\begin{itemize}
\tightlist
\item
  Has \(\chi_{\mathrm{cpl}} > 6\) (excluded by MAR), or
\item
  Fails to provide both pseudoreal and complex representations (excluded by condition iii), or
\item
  Is a subgroup of the SM connected group, which would fail anomaly cancellation or Yukawa completeness.
\end{itemize}

\(\square\)

\subsubsection{Corollary 3: Hypercharge quantization}\label{corollary-3-hypercharge-quantization}

\textbf{Corollary.} The hypercharge ratios are fixed by anomaly cancellation and Yukawa invariance within the derived gauge group, and the standard normalization then fixes the observed values \(Y \in \frac{1}{6}\mathbb{Z}\).

\emph{Proof.} The hypercharge theorem in the main OPH trunk derives the unique anomaly-free hypercharge ratios for \(SU(3) \times SU(2) \times U(1)\) with \(N_c = 3\), assuming single-Higgs Yukawa completeness (condition iv); the standard normalization then gives the sixth-integer lattice. \(\square\)

\begin{center}\rule{0.5\linewidth}{0.5pt}\end{center}

\subsection{Complete Derivation Chain}\label{9-complete-derivation-chain}

The gauge-sector chain established in this fragment is:

\begin{quote}
\(\text{canonical package} + \text{fixed-cutoff realized presentation}\) \(\xrightarrow{\text{Thm 6.1}}\) \(\exists\, G\) compact

\(+ \; [z]=0\) \(\xrightarrow{\text{Prop 6.1a}}\) DHR transportable, global gauge symmetry

\(+ \;\text{apply MAR on the admissible class}\) \(\xrightarrow{\text{Thm (this doc)}}\) connected SM factors with the realized chain \(N_g = 3\) then \(N_c = 3\), and \(G_{\mathrm{phys}} = SU(3) \times SU(2) \times U(1)/\mathbb{Z}_6\)

\(\xrightarrow{\text{Thm 6.13}}\) hypercharges fixed \(\xrightarrow{\text{Cor.}}\) no gauge-mediated proton decay
\end{quote}

Under the gauge-selection hypotheses above, the following are all derived results:

\begin{enumerate}
\def\labelenumi{\arabic{enumi}.}
\tightlist
\item
  Product gauge group structure
\item
  Exact global gauge group \(SU(3) \times SU(2) \times U(1)/\mathbb{Z}_6\)
\item
  Number of colors \(N_c = 3\)
\item
  Number of generations \(N_g = 3\)
\item
  Absence of gauge-mediated proton decay
\item
  Hypercharge quantization in sixths
\end{enumerate}

This is the gauge-sector closure on the stated admissible class. Gravity, spectrum, and other phenomenological branches use additional premises stated in their own sections. \(\square\)
